Knowledge graphs are republished on a schedule, and every release adds, removes
and rewrites facts. Between YAGO~4 and YAGO~4.5.0.2 the class
\texttt{schema:Person} lost 245{,}398 triples and gained 223{,}227 inside a
window of 65{,}182 subjects, while keeping its name and its identifier.
Consumers inherit that change without being told. Plenty of metrics describe one
release; almost nothing describes the difference between two at the scale the
releases actually have.

This thesis builds a catalogue of thirteen metrics spanning the class, entity
and global level, implemented in Rust against any QLever SPARQL endpoint, and a
dashboard that places the results of several releases side by side. Two
components are contributions rather than adoptions. Entropy-weighted PageRank
weights every edge by the normalised property entropy of its predicate, so score
flows through informative relations rather than high-volume bookkeeping ones. An
integer-encoded triple diff with subject-window scoping makes an exact version
comparison possible without a global sort, which SPARQL engines refuse on graphs
of this size.

The evaluation covers seven snapshots: two generations of YAGO, a decade of
DBpedia, and the pre-schema.org YAGO~3 release. The integer encoding runs 2.34 to
6.17 times faster than a string-based baseline and holds its working set in 16 to
19 times less memory, with both paths asserted to return identical results. The
memory figure is what decides feasibility on a 16\,GB client. The metrics show
the two graphs evolving in opposite directions: YAGO replaced 94\% of its class
vocabulary while keeping 72 predicates, whereas DBpedia kept 700 of 740 classes
while 30{,}788 predicates disappeared. Across graphs, DBpedia uses 24 to 92
times more predicates per class than YAGO on all sixteen classes matched
automatically between the latest releases. \texttt{Organization} is the only
class where YAGO~4.5.0.2 and DBpedia 2025 report the same entropy, 14.275 bits
on each side, despite differing in size by a factor of 3.7. Repeating the
comparison against four DBpedia releases and against YAGO~4 shows the agreement
is recent and needed both graphs to change.
